% Options for packages loaded elsewhere
\PassOptionsToPackage{unicode}{hyperref}
\PassOptionsToPackage{hyphens}{url}
\PassOptionsToPackage{dvipsnames,svgnames,x11names}{xcolor}
\documentclass[
  spanish,
  a4paper,
  12pt,
]{article}
\usepackage{xcolor}
\usepackage{graphicx}
\usepackage{tikz}
\usetikzlibrary{calc}
\usepackage{float}
\usepackage[margin=2.5cm]{geometry}
\usepackage{amsmath,amssymb}
\setcounter{secnumdepth}{5}
\linespread{1.12}
\usepackage{iftex}
\ifPDFTeX
  \usepackage[T1]{fontenc}
  \usepackage[utf8]{inputenc}
  \usepackage{textcomp} % provide euro and other symbols
\else % if luatex or xetex
  \usepackage{unicode-math} % this also loads fontspec
  \defaultfontfeatures{Scale=MatchLowercase}
  \defaultfontfeatures[\rmfamily]{Ligatures=TeX,Scale=1}
\fi
\usepackage{lmodern}
\ifPDFTeX\else
  % xetex/luatex font selection
\fi
% Use upquote if available, for straight quotes in verbatim environments
\IfFileExists{upquote.sty}{\usepackage{upquote}}{}
\IfFileExists{microtype.sty}{% use microtype if available
  \usepackage[]{microtype}
  \UseMicrotypeSet[protrusion]{basicmath} % disable protrusion for tt fonts
}{}
\makeatletter
\@ifundefined{KOMAClassName}{% if non-KOMA class
  \IfFileExists{parskip.sty}{%
    \usepackage{parskip}
  }{% else
    \setlength{\parindent}{0pt}
    \setlength{\parskip}{6pt plus 2pt minus 1pt}}
}{% if KOMA class
  \KOMAoptions{parskip=half}}
\makeatother
\usepackage{listings}
\usepackage{inconsolata}
\newcommand{\passthrough}[1]{#1}
\lstdefinestyle{codigo}{
  basicstyle=\ttfamily\small,
  breaklines=true,
  breakatwhitespace=false,
  columns=fullflexible,
  keepspaces=true,
  showstringspaces=false,
  frame=single,
  rulecolor=\color{gray!35},
  backgroundcolor=\color{gray!6},
  xleftmargin=0.5em,
  xrightmargin=0.5em,
  aboveskip=0.8em,
  belowskip=0.8em,
  inputencoding=utf8,
  extendedchars=true,
  literate=
    {á}{{\'a}}1 {é}{{\'e}}1 {í}{{\'i}}1 {ó}{{\'o}}1 {ú}{{\'u}}1
    {Á}{{\'A}}1 {É}{{\'E}}1 {Í}{{\'I}}1 {Ó}{{\'O}}1 {Ú}{{\'U}}1
    {ñ}{{\~n}}1 {Ñ}{{\~N}}1
    {°}{{$^\circ$}}1 {≤}{{$\leq$}}1 {≥}{{$\geq$}}1
    {→}{{$\to$}}1 {←}{{$\leftarrow$}}1
    {—}{{---}}1 {–}{{--}}1 {−}{{$-$}}1
    {×}{{$\times$}}1 {≈}{{$\approx$}}1
    {├}{{|}}1 {└}{{\textbackslash}}1 {│}{{|}}1 {─}{{-}}1
}
\lstset{style=codigo}
\usepackage{longtable,booktabs,array}
\newcounter{none} % for unnumbered tables
\usepackage{calc} % for calculating minipage widths
% Correct order of tables after \paragraph or \subparagraph
\usepackage{etoolbox}
\makeatletter
\patchcmd\longtable{\par}{\if@noskipsec\mbox{}\fi\par}{}{}
\makeatother
% Allow footnotes in longtable head/foot
\IfFileExists{footnotehyper.sty}{\usepackage{footnotehyper}}{\usepackage{footnote}}
\makesavenoteenv{longtable}
\ifLuaTeX
\usepackage[bidi=basic,shorthands=off]{babel}
\else
\usepackage[bidi=default,shorthands=off]{babel}
\fi
\ifLuaTeX
  \usepackage{selnolig} % disable illegal ligatures
\fi
\setlength{\emergencystretch}{3em} % prevent overfull lines
\providecommand{\tightlist}{%
  \setlength{\itemsep}{0pt}\setlength{\parskip}{0pt}}
\usepackage{bookmark}
\IfFileExists{xurl.sty}{\usepackage{xurl}}{} % add URL line breaks if available
\urlstyle{same}
\definecolor{negro_suave}{rgb}{.15,.15,.15}
\hypersetup{
  pdflang={es},
  colorlinks=true,
  linkcolor={blue},
  filecolor={Maroon},
  citecolor={Blue},
  urlcolor={blue},
  pdftitle={Informe Técnico: Algoritmos de Rasterización y Relleno en Computación Gráfica},
  pdfauthor={Julio Cesar Blacio Machuca},
  pdfcreator={LaTeX via pandoc}}

% Logos de portada: sube estos archivos a Overleaf con esos nombres.
\newcommand{\LogoUniversidad}{Logo.png}
\newcommand{\LogoDepartamento}{DCCO.png}
\newcommand{\InsertLogo}[2]{%
  \IfFileExists{#1}{\includegraphics[height=#2]{#1}}{%
    \fbox{\parbox[c][#2][c]{#2}{\centering\scriptsize Logo}}%
  }%
}
\newcommand{\InsertLogoWidth}[2]{%
  \IfFileExists{#1}{\includegraphics[width=#2]{#1}}{%
    \fbox{\parbox[c][1.4cm][c]{#2}{\centering\scriptsize Logo}}%
  }%
}
\newcommand{\CapturaPrueba}[3]{%
  \begin{figure}[H]
  \centering
  \IfFileExists{#1}{%
    \includegraphics[width=0.9\linewidth]{#1}%
  }{%
    \fbox{\parbox[c][6cm][c]{0.9\linewidth}{%
      \centering
      \textbf{Espacio reservado para captura}\\[0.4cm]
      Archivo esperado: \texttt{#1}\\[0.2cm]
      #3
    }}%
  }%
  \caption{#2}
  \end{figure}
}

\title{Informe Técnico: Algoritmos de Rasterización y Relleno en Computación Gráfica}
\author{Julio Cesar Blacio Machuca}
\date{Junio 2026}

\begin{document}

\begin{titlepage}
\thispagestyle{empty}
\begin{tikzpicture}[remember picture,overlay]
  % Barras decorativas como en Alter.tex
  \filldraw[fill=negro_suave, draw=negro_suave]
    ($(current page.north west)+(2.0cm,-2.7cm)$)
    rectangle
    ($(current page.south west)+(2.35cm,2.0cm)$);
  \filldraw[fill=negro_suave, draw=negro_suave]
    ($(current page.north west)+(3.2cm,-4.05cm)$)
    rectangle
    ($(current page.north east)+(-2.1cm,-3.85cm)$);
  \filldraw[fill=negro_suave, draw=negro_suave]
    ($(current page.south west)+(3.2cm,3.3cm)$)
    rectangle
    ($(current page.south east)+(-2.1cm,3.5cm)$);

  % Logos. Sube a Overleaf Logo.png y DCCO.png.
  \node[anchor=north west] at ($(current page.north west)+(3.35cm,-1.95cm)$)
    {\InsertLogoWidth{\LogoUniversidad}{6.0cm}};
  \node[anchor=north east] at ($(current page.north east)+(-2.55cm,-1.65cm)$)
    {\InsertLogo{\LogoDepartamento}{2.6cm}};

  % Texto principal de portada
  \node[anchor=east] at ($(current page.east)+(-3.0cm,1.4cm)$)
    {\textbf{\Huge Computación Gráfica}};
  \node[anchor=east] at ($(current page.east)+(-3.0cm,-0.3cm)$)
    {\textbf{\LARGE INFORME DEL PROYECTO}};
  \node[anchor=east] at ($(current page.east)+(-3.0cm,-2.2cm)$)
    {\Large Autor: Julio Cesar Blacio Machuca};
  \node[anchor=east] at ($(current page.east)+(-3.0cm,-3.4cm)$)
    {\Large NRC: 29505};
  \node[anchor=east] at ($(current page.east)+(-3.0cm,-4.6cm)$)
    {\Large Fecha: Junio 2026};
  \node[anchor=east] at ($(current page.south east)+(-2.9cm,4.45cm)$)
    {\textbf{\Huge ALGORITMOS RÁSTER}};

  % Pie de portada
  \node[anchor=east] at ($(current page.south east)+(-2.6cm,1.9cm)$)
    {\textit{{\Large ING. DARIO JAVIER MORALES CAIZA}}};
  \node[anchor=east] at ($(current page.south east)+(-2.6cm,1.15cm)$)
    {\textit{{\Large Departamento de Ciencias de la Computación \quad ESPE}}};
\end{tikzpicture}
\mbox{}
\end{titlepage}

{
\setcounter{tocdepth}{3}
\tableofcontents
}
\newpage

\section{Resumen}\label{resumen}

El presente informe documenta el diseño, implementación y validación de
nueve algoritmos clásicos de computación gráfica interactiva,
desarrollados sobre la plataforma .NET 10.0 en C\# con Windows Forms. El
sistema está organizado bajo el patrón arquitectónico
Modelo-Vista-Controlador (MVC) y abarca tres categorías de algoritmos:
rasterización de líneas rectas (DDA, Bresenham, Punto Medio de Línea),
generación de circunferencias (Punto Medio de Círculo, Paramétrico,
Polar) y algoritmos de relleno de regiones (por Semilla/BFS, por
Pila/DFS, Scanline). Cada algoritmo cuenta con una interfaz gráfica
interactiva que permite visualización paso a paso, control de parámetros
en tiempo real y animación con retardo configurable. Se documentan las
decisiones de diseño de interfaz, el sistema de coordenadas, el manejo
de excepciones y la complejidad computacional comparativa de cada
enfoque.

\begin{center}\rule{0.5\linewidth}{0.5pt}\end{center}

\section{Arquitectura del Sistema y Decisiones de Diseño
UI/UX}\label{arquitectura-del-sistema-y-decisiones-de-diseuxf1o-uiux}

\subsection{Patrón Arquitectónico
MVC}\label{patruxf3n-arquitectuxf3nico-mvc}

El sistema implementa el patrón \textbf{Modelo-Vista-Controlador} de
forma estricta a nivel de directorios y namespaces. Esta separación
garantiza que la lógica algorítmica sea completamente independiente de
la capa de presentación, facilitando el mantenimiento y la
extensibilidad.

\begin{lstlisting}
GeometriaComputacional/
├── Lineas/
│   ├── DDA/
│   │   ├── Controlador/   DdaController.cs
│   │   ├── Modelo/        PasoDDA.cs, ResultadoDDA.cs
│   │   └── Vista/         FrmDDA.cs
│   ├── Bresenham/         (misma estructura)
│   └── PuntoMedioLinea/   (misma estructura)
├── Conicas/
│   ├── PuntoMedioCirculo/ (misma estructura)
│   ├── ParametricoCirculo/
│   └── PolarCirculo/
└── Relleno/
    ├── RellenoPorSemilla/
    ├── RellenoPorPila/
    └── RellenoScanline/
\end{lstlisting}

Cada \textbf{Controlador} expone únicamente un método
\passthrough{\lstinline!Calcular(...)!} que recibe los parámetros
matemáticos y retorna un objeto \textbf{Resultado} inmutable. La
\textbf{Vista} consume ese resultado a través de un temporizador o de
forma directa, sin conocer los detalles algorítmicos.

\subsection{Paleta de Colores con Canal Alfa
(ARGB)}\label{paleta-de-colores-con-canal-alfa-argb}

Todos los elementos gráficos del sistema utilizan el modelo de color
\textbf{ARGB} de 32 bits, donde el canal alfa (primeros 8 bits) controla
la transparencia con valores entre 0 (totalmente transparente) y 255
(totalmente opaco).

La elección de transparencia parcial responde a una necesidad funcional
concreta: los \textbf{píxeles rasterizados} deben superponerse
visualmente sobre la grilla del plano cartesiano y la línea/curva ideal
sin ocultarla completamente, de modo que el usuario pueda observar
simultáneamente la aproximación discreta y la figura matemática
continua.

{\def\LTcaptype{none} % do not increment counter
\begin{longtable}[]{@{}
  >{\raggedright\arraybackslash}p{(\linewidth - 6\tabcolsep) * \real{0.2295}}
  >{\raggedright\arraybackslash}p{(\linewidth - 6\tabcolsep) * \real{0.2869}}
  >{\raggedright\arraybackslash}p{(\linewidth - 6\tabcolsep) * \real{0.0574}}
  >{\raggedright\arraybackslash}p{(\linewidth - 6\tabcolsep) * \real{0.4262}}@{}}
\toprule\noalign{}
\begin{minipage}[b]{\linewidth}\raggedright
Elemento
\end{minipage} & \begin{minipage}[b]{\linewidth}\raggedright
Color ARGB
\end{minipage} & \begin{minipage}[b]{\linewidth}\raggedright
Alpha
\end{minipage} & \begin{minipage}[b]{\linewidth}\raggedright
Justificación
\end{minipage} \\
\midrule\noalign{}
\endhead
\bottomrule\noalign{}
\endlastfoot
Píxeles DDA / Bresenham & \passthrough{\lstinline!(255, 220, 38, 38)!} &
255 & Rojo sólido, contraste máximo sobre plano blanco \\
Círculo ideal de referencia &
\passthrough{\lstinline!(160, 37, 99, 235)!} & 160 & Azul
semitransparente; visible bajo los píxeles \\
Píxeles paramétrico & \passthrough{\lstinline!(200, 22, 163, 74)!} & 200
& Verde con ligera transparencia \\
Píxeles polar & \passthrough{\lstinline!(200, 147, 51, 234)!} & 200 &
Morado semitransparente \\
Relleno BFS (semilla) & \passthrough{\lstinline!(255, 186, 230, 253)!} &
255 & Azul claro, evoca expansión uniforme \\
Relleno DFS (pila) & \passthrough{\lstinline!(255, 167, 243, 208)!} &
255 & Verde claro, evoca profundidad de árbol \\
Relleno Scanline & \passthrough{\lstinline!(255, 233, 213, 255)!} & 255
& Lila claro, evoca barrido horizontal \\
\end{longtable}
}

\subsection{Suavizado
Antialiasing}\label{suavizado-antialiasing}

Todos los formularios de líneas y cónicas activan el modo de suavizado
de alta calidad en el evento \passthrough{\lstinline!Paint!}:

\begin{lstlisting}
g.SmoothingMode = System.Drawing.Drawing2D.SmoothingMode.AntiAlias;
\end{lstlisting}

Esta instrucción indica a GDI+ que aplique \textbf{submuestreo
multimuestra} al renderizar primitivas vectoriales (líneas, elipses,
arcos). Sin esta activación, los bordes de las figuras ideales presentan
el efecto ``escalera'' (\emph{jaggies}), lo que dificulta la comparación
visual entre la figura matemática continua y los píxeles rasterizados.
El antialiasing se aplica únicamente a los elementos vectoriales; los
rectángulos de píxeles rasterizados se dibujan intencionalmente sin
suavizado para representar fielmente la naturaleza discreta del
resultado del algoritmo.

\subsection{Orden de Pintado (Z-order de
Software)}\label{orden-de-pintado-z-order-de-software}

Una decisión crítica de renderizado es el \textbf{orden de pintado} de
los elementos. GDI+ no gestiona capas; el último elemento dibujado
aparece al frente. El orden correcto garantiza que los puntos de control
siempre sean visibles:

\begin{lstlisting}
1. Píxeles rasterizados      ← fondo (pueden ser tapados)
2. Línea / curva ideal       ← nivel intermedio
3. Puntos de control (P1/P2) ← siempre al frente
\end{lstlisting}

Sin este orden, los cuadrados de píxeles rojos cubrirían los puntos de
control verdes y naranjas cuando un píxel coincide con la posición de un
punto de control, haciéndolos inaccesibles para el arrastre con el
ratón.

\begin{center}\rule{0.5\linewidth}{0.5pt}\end{center}

\section{Clasificación Taxonómica de
Algoritmos}\label{clasificaciuxf3n-taxonuxf3mica-de-algoritmos}

\begin{lstlisting}
Algoritmos de Computación Gráfica Implementados
│
├── RASTERIZACIÓN DE PRIMITIVAS LINEALES
│   ├── DDA (Digital Differential Analyzer)     — Incremental con aritmética real
│   ├── Bresenham                                — Incremental con aritmética entera
│   └── Punto Medio de Línea                    — Decisión por evaluación de función implícita
│
├── GENERACIÓN DE CIRCUNFERENCIAS
│   ├── Punto Medio de Círculo                  — Decisión incremental + simetría 8-fold
│   ├── Paramétrico                              — Muestreo trigonométrico uniforme por ángulo
│   └── Polar                                    — Conversión de coordenadas polares a cartesianas
│
└── RELLENO DE REGIONES CERRADAS
    ├── Por Semilla (BFS/Cola)                  — Expansión uniforme por capas
    ├── Por Pila (DFS/Stack)                    — Exploración en profundidad
    └── Scanline                                 — Relleno por tramos horizontales
\end{lstlisting}

\begin{center}\rule{0.5\linewidth}{0.5pt}\end{center}

\section{Capturas de Pantalla de Pruebas Realizadas}\label{capturas-de-pantalla-de-pruebas-realizadas}

Esta sección reúne las evidencias visuales más importantes del
funcionamiento del sistema. Las capturas seleccionadas cubren las tres
familias de algoritmos implementadas: rasterización de líneas,
generación de circunferencias y relleno de regiones. Se recomienda usar
capturas donde se observe simultáneamente el plano o matriz de trabajo,
los parámetros ingresados, la tabla de pasos y el resultado gráfico.

\subsection{Prueba de Rasterización de Líneas}\label{prueba-de-rasterizacion-de-lineas}

La captura de líneas debe mostrar la comparación entre la recta ideal y
los píxeles discretos generados por el algoritmo. La prueba más
representativa es Bresenham o DDA, porque permite observar claramente la
diferencia entre una línea continua y su aproximación rasterizada sobre
la grilla.

\CapturaPrueba
  {captura-linea-bresenham.png}
  {Prueba del algoritmo de Bresenham para rasterización de líneas.}
  {Sugerencia: captura la ventana con los puntos inicial y final, la línea ideal, los píxeles generados y la tabla de pasos.}

\subsection{Prueba de Generación de Circunferencias}\label{prueba-de-generacion-de-circunferencias}

La captura de circunferencias debe evidenciar la simetría del algoritmo
respecto al centro y la forma en que los puntos generados aproximan la
circunferencia ideal. La prueba más útil es Punto Medio de Círculo,
porque muestra la ventaja de calcular un octante y reflejar los puntos
por simetría.

\CapturaPrueba
  {captura-circulo-punto-medio.png}
  {Prueba del algoritmo de Punto Medio para generación de circunferencias.}
  {Sugerencia: captura el centro, el radio, la circunferencia ideal, los puntos rasterizados y la tabla de iteraciones.}

\subsection{Prueba de Relleno de Regiones}\label{prueba-de-relleno-de-regiones}

La captura de relleno debe mostrar una región cerrada, la semilla
seleccionada y el patrón de expansión del algoritmo. La prueba más
representativa es Scanline, porque permite comparar visualmente su
eficiencia frente a BFS y DFS al rellenar tramos horizontales completos.

\CapturaPrueba
  {captura-relleno-scanline.png}
  {Prueba del algoritmo Scanline para relleno de regiones.}
  {Sugerencia: captura la figura cerrada, la semilla, las celdas rellenas y el avance por filas.}



\begin{center}\rule{0.5\linewidth}{0.5pt}\end{center}

\section{Algoritmos de Rasterización de
Líneas}\label{algoritmos-de-rasterizaciuxf3n-de-luxedneas}

\subsection{Algoritmo DDA (Digital Differential
Analyzer)}\label{algoritmo-dda-digital-differential-analyzer}

\subsubsection{Análisis Matemático}\label{anuxe1lisis-matemuxe1tico}

El DDA convierte una recta continua \(f(x) = mx + b\) en una secuencia
de píxeles discretos mediante incrementos acumulativos. Dados dos puntos
\(P_1 = (x_1, y_1)\) y \(P_2 = (x_2, y_2)\):

\[\Delta x = x_2 - x_1 \qquad \Delta y = y_2 - y_1\]

El número de pasos \(k\) determina cuántos píxeles se generan:

\[k = \max(|\Delta x|,\ |\Delta y|)\]

Los incrementos por paso son fracciones racionales:

\[\delta x = \frac{\Delta x}{k} \qquad \delta y = \frac{\Delta y}{k}\]

La secuencia de píxeles se calcula iterativamente:

\[x_{i+1} = x_i + \delta x \qquad y_{i+1} = y_i + \delta y\]

El píxel en la pantalla se obtiene redondeando:

\[P_i = \bigl(\text{round}(x_i),\ \text{round}(y_i)\bigr)\]

\textbf{Propiedad clave:} El incremento mayor siempre es \(\pm 1\),
garantizando que no haya píxeles saltados y que la recta sea conexa.

\newpage

\subsubsection{Fragmento de Código}\label{fragmento-de-cuxf3digo}

\begin{lstlisting}
// DdaController.cs — método Calcular
public ResultadoDDA Calcular(PuntoGeometrico p1, PuntoGeometrico p2)
{
    float dx = p2.X - p1.X;
    float dy = p2.Y - p1.Y;

    // k = número total de píxeles a generar
    float k = Math.Max(Math.Abs(dx), Math.Abs(dy));

    // Caso degenerado: ambos puntos coinciden
    if (k == 0)
        return new ResultadoDDA(0, 0, 0, 1, new List<PasoDDA>
            { new(0, p1.X, p1.Y, (int)Math.Round(p1.X), (int)Math.Round(p1.Y)) });

    float ix = dx / k;   // incremento fraccional en X
    float iy = dy / k;   // incremento fraccional en Y

    float x = p1.X, y = p1.Y;
    var pasos = new List<PasoDDA>();

    for (int i = 0; i <= (int)k; i++)
    {
        // Redondeo al píxel más cercano
        int px = (int)Math.Round(x);
        int py = (int)Math.Round(y);
        pasos.Add(new PasoDDA(i, x, y, px, py));
        x += ix;
        y += iy;
    }

    float pendiente = dx == 0 ? float.PositiveInfinity : dy / dx;
    return new ResultadoDDA(dx, dy, pendiente, (int)k + 1, pasos);
}
\end{lstlisting}

\newpage

\subsubsection{Resultado de Prueba}\label{resultado-de-prueba}

{\def\LTcaptype{none} % do not increment counter
\begin{longtable}[]{@{}ll@{}}
\toprule\noalign{}
Parámetros & Valores \\
\midrule\noalign{}
\endhead
\bottomrule\noalign{}
\endlastfoot
\(P_1\) & \((-12,\ -2)\) \\
\(P_2\) & \((11,\ -1)\) \\
\(\Delta x\) & \(23\) \\
\(\Delta y\) & \(1\) \\
\(k\) & \(23\) pasos \\
\(\delta x\) & \(1.000\) \\
\(\delta y\) & \(0.043\) \\
Pendiente \(m\) & \(0.04\) \\
\end{longtable}
}

Los píxeles generados pasan de \((-12, -2)\) en pasos unitarios en
\(x\). Dado que \(\delta y \approx 0.043\), la coordenada \(y\) solo
cambia de \(-2\) a \(-1\) una vez en el paso 12, reflejando fielmente la
pendiente casi nula de la recta.

\begin{center}\rule{0.5\linewidth}{0.5pt}\end{center}

\subsection{Algoritmo de
Bresenham}\label{algoritmo-de-bresenham}

\subsubsection{Análisis
Matemático}\label{anuxe1lisis-matemuxe1tico-1}

Bresenham elimina el uso de aritmética de punto flotante al trabajar
exclusivamente con \textbf{enteros} y una \textbf{variable de decisión}
\(e\) que acumula el error respecto a la recta ideal.

Para el caso general con \(|\Delta x| \geq |\Delta y|\) (octante 1):

\[e_0 = 2|\Delta y| - |\Delta x|\]

En cada iteración se evalúa \(e_2 = 2e\):

\[\text{Si } e_2 > -|\Delta y|: \quad x \mathrel{+}= s_x, \quad e \mathrel{-}= 2|\Delta y|\]
\[\text{Si } e_2 < |\Delta x|: \quad y \mathrel{+}= s_y, \quad e \mathrel{+}= 2|\Delta x|\]

Donde \(s_x = \text{signo}(\Delta x)\) y
\(s_y = \text{signo}(\Delta y)\) son las direcciones de avance. Este
criterio generalizado cubre todos los octantes sin casos especiales.

\newpage

\subsubsection{Fragmento de Código}\label{fragmento-de-cuxf3digo-1}

\begin{lstlisting}
// BresenhamController.cs — algoritmo generalizado para todos los octantes
int x = p1.X, y = p1.Y;
int dx = Math.Abs(p2.X - p1.X), dy = Math.Abs(p2.Y - p1.Y);
int sx = p1.X < p2.X ? 1 : -1;   // dirección en X
int sy = p1.Y < p2.Y ? 1 : -1;   // dirección en Y
int error = dx - dy;               // variable de decisión inicial

while (true)
{
    pasos.Add(new PasoBresenham(indice++, x, y, error));

    if (x == p2.X && y == p2.Y) break;

    int e2 = 2 * error;

    // Decide avance en X
    if (e2 > -dy) { error -= dy; x += sx; }

    // Decide avance en Y
    if (e2 < dx)  { error += dx; y += sy; }
}
\end{lstlisting}


\subsubsection{Resultado de Prueba}\label{resultado-de-prueba-1}

{\def\LTcaptype{none} % do not increment counter
\begin{longtable}[]{@{}lllll@{}}
\toprule\noalign{}
Paso & \(x\) & \(y\) & Error \(e\) & Acción \\
\midrule\noalign{}
\endhead
\bottomrule\noalign{}
\endlastfoot
0 & 0 & 0 & 3 & Inicio \\
1 & 1 & 0 & 1 & Avance en X \\
2 & 2 & 0 & −1 & Avance en X \\
3 & 2 & 1 & 5 & Avance en X e Y \\
4 & 3 & 1 & 3 & Avance en X \\
\end{longtable}
}

Para la recta de \(P_1=(0,0)\) a \(P_2=(5,3)\): \(\Delta x=5\),
\(\Delta y=3\), \(e_0=2\). El algoritmo produce 6 píxeles sin ninguna
operación de punto flotante.

\begin{center}\rule{0.5\linewidth}{0.5pt}\end{center}

\subsection{Algoritmo de Punto Medio de
Línea}\label{algoritmo-de-punto-medio-de-luxednea}

\subsubsection{Análisis
Matemático}\label{anuxe1lisis-matemuxe1tico-2}

El algoritmo de Punto Medio evalúa la función implícita de la recta
\(F(x,y) = \Delta y \cdot x - \Delta x \cdot y + c\) en el \textbf{punto
medio} entre los dos candidatos de píxel para decidir cuál está más
cerca de la recta ideal.

Parámetros derivados (asumiendo \(|\Delta x| \geq |\Delta y|\),
\(m \leq 1\)):

\[\Delta_m = \min(|\Delta x|,\ |\Delta y|), \qquad \Delta_M = \max(|\Delta x|,\ |\Delta y|)\]

Variable de decisión inicial:

\[p_0 = 2\Delta_m - \Delta_M\]

Actualización en cada paso:

\[\text{Si } p < 0: \quad p \mathrel{+}= 2\Delta_m \qquad \text{(avance recto)}\]
\[\text{Si } p \geq 0: \quad p \mathrel{+}= 2(\Delta_m - \Delta_M) \qquad \text{(avance diagonal)}\]

La \textbf{decisión} registrada en la tabla
(\passthrough{\lstinline!Recto!} o \passthrough{\lstinline!Diagonal!})
muestra directamente qué candidato de píxel fue elegido en cada
iteración.

\newpage

\subsubsection{Fragmento de Código}\label{fragmento-de-cuxf3digo-2}

\begin{lstlisting}
// PuntoMedioLineaController.cs
int deltaM = Math.Max(Math.Abs(dx), Math.Abs(dy));  // mayor delta
int deltam = Math.Min(Math.Abs(dx), Math.Abs(dy));  // menor delta

int p = 2 * deltam - deltaM;   // decisión inicial

for (int i = 0; i < deltaM; i++)
{
    string decision;
    if (p < 0)
    {
        // El punto medio indica: avanzar recto (solo eje dominante)
        p += 2 * deltam;
        decision = "Recto";
    }
    else
    {
        // El punto medio indica: avanzar diagonal
        p += 2 * (deltam - deltaM);
        decision = "Diagonal";
        avanzarEjeSecundario();
    }
    avanzarEjeDominante();
    pasos.Add(new PasoPuntoMedioLinea(i, xReal, yReal, xPixel, yPixel, decision));
}
\end{lstlisting}

\subsubsection{Resultado de Prueba}\label{resultado-de-prueba-2}

Para \(P_1=(0,0)\), \(P_2=(8,3)\): \(\Delta x=8\), \(\Delta y=3\),
\(p_0 = 2(3)-8 = -2\).

{\def\LTcaptype{none} % do not increment counter
\begin{longtable}[]{@{}llll@{}}
\toprule\noalign{}
Paso & Píxel & \(p\) & Decisión \\
\midrule\noalign{}
\endhead
\bottomrule\noalign{}
\endlastfoot
0 & \((0,0)\) & −2 & Recto \\
1 & \((1,0)\) & 4 & Recto \\
2 & \((2,1)\) & −8 & Diagonal \\
3 & \((3,1)\) & −2 & Recto \\
4 & \((4,1)\) & 4 & Recto \\
5 & \((5,2)\) & −8 & Diagonal \\
6 & \((6,2)\) & −2 & Recto \\
7 & \((7,2)\) & 4 & Recto \\
\end{longtable}
}

\begin{center}\rule{0.5\linewidth}{0.5pt}\end{center}

\section{Algoritmos de Generación de
Circunferencias}\label{algoritmos-de-generaciuxf3n-de-circunferencias}

\subsection{Punto Medio de
Círculo}\label{punto-medio-de-cuxedrculo}

\subsubsection{Análisis
Matemático}\label{anuxe1lisis-matemuxe1tico-3}

El algoritmo de Punto Medio para circunferencias explota la
\textbf{simetría de 8 octantes}: calcula los píxeles en el primer
octante (\(0° \leq \theta \leq 45°\), zona donde \(x \leq y\)) y los
refleja simétricamente.

La función implícita de la circunferencia es:

\[f(x, y) = x^2 + y^2 - r^2\]

Condición inicial:

\[x_0 = 0, \quad y_0 = r, \quad p_0 = 1 - r\]

Reglas de actualización en cada iteración (\(x \mathrel{+}= 1\)
siempre):

\[\text{Si } p < 0: \qquad p \mathrel{+}= 2x + 3\]

\[\text{Si } p \geq 0: \qquad y \mathrel{-}= 1, \quad p \mathrel{+}= 2(x - y) + 5\]

\begin{quote}
\textbf{Nota de implementación crítica:} el decremento
\(y \mathrel{-}= 1\) se aplica \textbf{antes} de evaluar la fórmula
\(p + 2(x-y)+5\). El término \((x-y)\) utiliza el valor actualizado de
\(y\).
\end{quote}

Los 8 puntos simétricos generados a partir de \((x, y)\) relativo al
centro \((x_c, y_c)\) son:

\[(\pm x + x_c,\ \pm y + y_c) \qquad (\pm y + x_c,\ \pm x + y_c)\]

\newpage

\subsubsection{Fragmento de Código}\label{fragmento-de-cuxf3digo-3}

\begin{lstlisting}
// PuntoMedioCirculoController.cs
int x = 0, y = radio, p = 1 - radio, iter = 0;

while (x <= y)
{
    // Obtener los 8 puntos simétricos (con deduplicación por HashSet)
    var puntos = ObtenerSimetrias(xc, yc, x, y);
    pasos.Add(new PasoPuntoMedioCirculo(iter, x, y, p,
        FormatearSimetrias(puntos), puntos));

    // Actualizar variable de decisión
    if (p < 0)
    {
        p += 2 * x + 3;        // avance Este
    }
    else
    {
        y--;                   // avance Sureste: decrementar y primero
        p += 2 * (x - y) + 5; // luego evaluar con y ya decrementado
    }
    x++;
    iter++;
}
\end{lstlisting}

\subsubsection{Resultado de Prueba}\label{resultado-de-prueba-3}

Para \(x_c=0\), \(y_c=0\), \(r=5\):

{\def\LTcaptype{none} % do not increment counter
\begin{longtable}[]{@{}lllll@{}}
\toprule\noalign{}
Iter & \(x\) & \(y\) & \(p\) & Decisión \\
\midrule\noalign{}
\endhead
\bottomrule\noalign{}
\endlastfoot
0 & 0 & 5 & −4 & Este \\
1 & 1 & 5 & −2 & Este \\
2 & 2 & 5 & 2 & Sureste \\
3 & 3 & 4 & −6 & Este \\
4 & 4 & 4 & 2 & Sureste \\
\end{longtable}
}

En la iteración 2, \(p = -2 + 2(2)+3 = 5 \geq 0\), por lo que \(y\) pasa
de 5 a 4 y \(p = 5 + 2(2-4)+5 = 2\).

\begin{center}\rule{0.5\linewidth}{0.5pt}\end{center}

\subsection{Paramétrico de
Círculo}\label{paramuxe9trico-de-cuxedrculo}

\subsubsection{Análisis
Matemático}\label{anuxe1lisis-matemuxe1tico-4}

La representación paramétrica expresa cada punto de la circunferencia en
función del ángulo \(t\):

\[x(t) = x_c + r\cos(t) \qquad y(t) = y_c + r\sin(t)\]

donde \(t \in [0°, 360°)\) se muestrea con paso de \(1°\), dando
exactamente 360 puntos. El píxel correspondiente se obtiene redondeando
al entero más cercano:

\[P_t = \bigl(\text{round}(x(t)),\ \text{round}(y(t))\bigr)\]

El \textbf{paso angular de 1°} garantiza cobertura completa sin brechas
visibles para radios \(r \geq 2\), ya que el arco por grado es
\(s = \frac{2\pi r}{360} \approx 0.017r\), que para \(r=5\) da
\(s \approx 0.087 < 1\) px, asegurando que puntos consecutivos sean
adyacentes o coincidentes.

El total de píxeles únicos es siempre menor que 360 debido al redondeo:
múltiples ángulos pueden mapearse al mismo píxel, especialmente en los
ejes y las diagonales.

\subsubsection{Fragmento de Código}\label{fragmento-de-cuxf3digo-4}

\begin{lstlisting}
// ParametricoCirculoController.cs
var pasos = new List<PasoParametricoCirculo>();
var vistos = new HashSet<(int, int)>();

for (int t = 0; t < 360; t++)
{
    double rad = t * Math.PI / 180.0;
    float xR = (float)(xc + radio * Math.Cos(rad));
    float yR = (float)(yc + radio * Math.Sin(rad));

    // Redondeo al píxel más cercano
    int px = (int)Math.Round(xR);
    int py = (int)Math.Round(yR);

    pasos.Add(new PasoParametricoCirculo(t, xR, yR, px, py));

    // Conteo de píxeles únicos (sin duplicados por redondeo)
    vistos.Add((px, py));
}

return new ResultadoParametricoCirculo(xc, yc, radio, vistos.Count, pasos);
\end{lstlisting}

\subsubsection{Resultado de Prueba}\label{resultado-de-prueba-4}

Para \(x_c=0\), \(y_c=0\), \(r=5\):

{\def\LTcaptype{none} % do not increment counter
\begin{longtable}[]{@{}llll@{}}
\toprule\noalign{}
\(t\) (°) & \(x(t)\) real & \(y(t)\) real & Píxel \((p_x, p_y)\) \\
\midrule\noalign{}
\endhead
\bottomrule\noalign{}
\endlastfoot
0 & 5.000 & 0.000 & \((5, 0)\) \\
30 & 4.330 & 2.500 & \((4, 3)\) \\
45 & 3.536 & 3.536 & \((4, 4)\) \\
90 & 0.000 & 5.000 & \((0, 5)\) \\
180 & −5.000 & 0.000 & \((−5, 0)\) \\
270 & 0.000 & −5.000 & \((0, −5)\) \\
\end{longtable}
}

Total de 360 ángulos muestreados → \textbf{31 píxeles únicos} (para
\(r=5\)), evidenciando la redundancia por redondeo.

\begin{center}\rule{0.5\linewidth}{0.5pt}\end{center}

\subsection{Polar de Círculo}\label{polar-de-cuxedrculo}

\subsubsection{Análisis
Matemático}\label{anuxe1lisis-matemuxe1tico-5}

El método polar parte de las coordenadas polares \((r, \theta)\), donde
el \textbf{radio \(r\) es constante} y el ángulo \(\theta\) varía. La
conversión al plano cartesiano se realiza mediante:

\[x = x_c + r\cos(\theta) \qquad y = y_c + r\sin(\theta)\]

La diferencia conceptual respecto al paramétrico es que en el método
polar el punto de partida es la \textbf{definición geométrica de radio
constante} en coordenadas polares, enfatizando que todos los puntos de
la circunferencia están a la misma distancia \(r\) del centro, mientras
que el método paramétrico enfatiza el recorrido angular.

Matemáticamente, ambos producen la misma secuencia de píxeles para el
mismo \((x_c, y_c, r)\).

\newpage

\subsubsection{Fragmento de Código}\label{fragmento-de-cuxf3digo-5}

\begin{lstlisting}
// PolarCirculoController.cs
for (int theta = 0; theta < 360; theta++)
{
    double rad = theta * Math.PI / 180.0;

    // Conversión polar → cartesiano con centro trasladado
    int px = (int)Math.Round(xc + radio * Math.Cos(rad));
    int py = (int)Math.Round(yc + radio * Math.Sin(rad));

    // Se registra: PixelX, R constante, Theta, punto formateado
    pasos.Add(new PasoPolarCirculo(theta, radio, px, py));
    vistos.Add((px, py));
}
\end{lstlisting}

\subsubsection{Resultado de Prueba}\label{resultado-de-prueba-5}

Para \(x_c=0\), \(y_c=0\), \(r=5\):

{\def\LTcaptype{none} % do not increment counter
\begin{longtable}[]{@{}lll@{}}
\toprule\noalign{}
\(\theta\) (°) & \(r\) & \((p_x, p_y)\) \\
\midrule\noalign{}
\endhead
\bottomrule\noalign{}
\endlastfoot
0 & 5 & \((5, 0)\) \\
60 & 5 & \((3, 4)\) \\
120 & 5 & \((−3, 4)\) \\
180 & 5 & \((−5, 0)\) \\
240 & 5 & \((−3, −4)\) \\
300 & 5 & \((3, −4)\) \\
\end{longtable}
}

El radio \(r=5\) aparece como valor constante en todas las filas de la
tabla, ilustrando la naturaleza polar del algoritmo.

\begin{center}\rule{0.5\linewidth}{0.5pt}\end{center}

\section{Algoritmos de Relleno de
Regiones}\label{algoritmos-de-relleno-de-regiones}

Los algoritmos de relleno operan sobre una \textbf{matriz de bytes} de
\(23 \times 15\) celdas. Cada celda tiene uno de tres estados:

{\def\LTcaptype{none} % do not increment counter
\begin{longtable}[]{@{}lll@{}}
\toprule\noalign{}
Valor & Estado & Color visualizado \\
\midrule\noalign{}
\endhead
\bottomrule\noalign{}
\endlastfoot
\passthrough{\lstinline!0!} & Vacía & Blanco \\
\passthrough{\lstinline!1!} & Borde & Gris oscuro \\
\passthrough{\lstinline!2!} & Rellena & Color del algoritmo \\
\end{longtable}
}

La \textbf{semilla} \((s_x, s_y)\) debe ubicarse en una celda vacía
dentro de la figura cerrada. El sistema incluye una \textbf{validación
previa} basada en BFS que verifica si la semilla puede alcanzar el borde
de la matriz; en caso afirmativo, la semilla está fuera de la figura y
el relleno se rechaza.

\subsection{Relleno por Semilla
(BFS)}\label{relleno-por-semilla-bfs}

\subsubsection{Análisis
Matemático}\label{anuxe1lisis-matemuxe1tico-6}

El algoritmo de Relleno por Semilla utiliza una \textbf{cola FIFO}
(First In, First Out), implementando una búsqueda en anchura (BFS ---
\emph{Breadth-First Search}). La expansión se produce en capas
concéntricas a partir de la semilla, análoga a las ondas de un punto en
un estanque.

El orden de procesamiento de vecinos es:

\[\text{Arriba} \to \text{Derecha} \to \text{Abajo} \to \text{Izquierda}\]

\[\{(x, y-1),\ (x+1, y),\ (x, y+1),\ (x-1, y)\}\]

Este orden, combinado con la disciplina FIFO de la cola, produce una
expansión visual uniforme: todas las celdas a distancia Manhattan \(d\)
de la semilla se pintan antes que cualquier celda a distancia \(d+1\).

La \textbf{complejidad temporal} es \(O(n)\) donde \(n\) es el número de
celdas interiores, ya que cada celda se encola y desencola exactamente
una vez.

\newpage

\subsubsection{Fragmento de Código}\label{fragmento-de-cuxf3digo-6}

\begin{lstlisting}
// RellenoPorSemillaController.cs — BFS con validación previa
public ResultadoRelleno Calcular(byte[,] matriz, int semX, int semY)
{
    // Validación 1: la semilla no debe estar en un borde
    // Validación 2: la semilla debe estar dentro de la figura cerrada
    if (matriz[semY, semX] != VACIA || !EsDentroDelaFigura(matriz, semX, semY))
        return new ResultadoRelleno(semX, semY, new List<PasoRelleno>());

    var cola = new Queue<(int X, int Y)>();
    var enCola = new bool[ROWS, COLS];
    cola.Enqueue((semX, semY));
    enCola[semY, semX] = true;
    int orden = 1;

    while (cola.Count > 0)
    {
        var (x, y) = cola.Dequeue();
        if (matriz[y, x] != VACIA) continue;

        matriz[y, x] = 2;  // marcar como rellena
        pasos.Add(new PasoRelleno(x, y, orden++));

        // Encolar vecinos en orden BFS: Arriba, Derecha, Abajo, Izquierda
        foreach (var (nx, ny) in new[] { (x, y-1), (x+1, y), (x, y+1), (x-1, y) })
        {
            if (nx >= 0 && nx < COLS && ny >= 0 && ny < ROWS
                && matriz[ny, nx] == VACIA && !enCola[ny, nx])
            {
                enCola[ny, nx] = true;
                cola.Enqueue((nx, ny));
            }
        }
    }
    return new ResultadoRelleno(semX, semY, pasos);
}
\end{lstlisting}

\newpage

\subsubsection{Resultado de Prueba}\label{resultado-de-prueba-6}

\textbf{Figura:} Cuadrado (borde en columnas 4--18, filas 2--12).
\textbf{Semilla:} \((11, 7)\).

{\def\LTcaptype{none} % do not increment counter
\begin{longtable}[]{@{}lll@{}}
\toprule\noalign{}
Orden & Celda & Descripción \\
\midrule\noalign{}
\endhead
\bottomrule\noalign{}
\endlastfoot
1 & \((11, 7)\) & Semilla inicial \\
2 & \((11, 6)\) & Arriba de la semilla \\
3 & \((12, 7)\) & Derecha de la semilla \\
4 & \((11, 8)\) & Abajo de la semilla \\
5 & \((10, 7)\) & Izquierda de la semilla \\
\end{longtable}
}

La animación con delay de 200ms evidencia el patrón de expansión
concéntrica. Total de celdas rellenas: \(15 \times 11 = 165\) (interior
del cuadrado).

\begin{center}\rule{0.5\linewidth}{0.5pt}\end{center}

\subsection{Relleno por Pila (DFS)}\label{relleno-por-pila-dfs}

\subsubsection{Análisis
Matemático}\label{anuxe1lisis-matemuxe1tico-7}

El Relleno por Pila utiliza una \textbf{pila LIFO} (Last In, First Out),
implementando una búsqueda en profundidad (DFS --- \emph{Depth-First
Search}). A diferencia del BFS, el algoritmo explora un camino hasta su
máxima profundidad antes de retroceder.

Los vecinos se \textbf{empujan} a la pila en el orden:

\[\text{Arriba},\ \text{Derecha},\ \text{Abajo},\ \text{Izquierda}\]

Pero debido a la disciplina LIFO, se \textbf{procesan} en orden inverso:

\[\text{Izquierda} \to \text{Abajo} \to \text{Derecha} \to \text{Arriba}\]

Esto produce un patrón visual característico: el relleno ``se mete'' por
una dirección y avanza lo más lejos posible antes de volver a explorar
otras direcciones, generando un patrón de expansión no uniforme que
visualmente recuerda al recorrido de un laberinto.

\newpage

\subsubsection{Fragmento de Código}\label{fragmento-de-cuxf3digo-7}

\begin{lstlisting}
// RellenoPorPilaController.cs — DFS con pila explícita
if (matriz[semY, semX] != VACIA || !EsDentroDelaFigura(matriz, semX, semY))
    return new ResultadoRelleno(semX, semY, new List<PasoRelleno>());

var pila = new Stack<(int X, int Y)>();
var visitado = new bool[ROWS, COLS];
pila.Push((semX, semY));

while (pila.Count > 0)
{
    var (x, y) = pila.Pop();

    if (x < 0 || x >= COLS || y < 0 || y >= ROWS) continue;
    if (matriz[y, x] != VACIA || visitado[y, x]) continue;

    visitado[y, x] = true;
    matriz[y, x] = 2;
    pasos.Add(new PasoRelleno(x, y, orden++));

    // Push en orden: Arriba, Derecha, Abajo, Izquierda
    // LIFO → saldrá Izquierda primero
    pila.Push((x,   y-1));  // Arriba    (sale último)
    pila.Push((x+1, y  ));  // Derecha
    pila.Push((x,   y+1));  // Abajo
    pila.Push((x-1, y  ));  // Izquierda (sale primero)
}
\end{lstlisting}

\subsubsection{Resultado de Prueba}\label{resultado-de-prueba-7}

\textbf{Figura:} Triángulo con vértices \((11,2)\), \((4,12)\),
\((18,12)\). \textbf{Semilla:} \((11, 7)\).

Los primeros pasos evidencian el comportamiento DFS: desde la semilla
\((11,7)\), el primer vecino procesado es Izquierda \((10,7)\). Desde
\((10,7)\) vuelve a tomar Izquierda, produciendo el recorrido:

\[S=(11,7) \to (10,7) \to (9,7) \to (8,7) \to \ldots\]

hasta alcanzar el borde izquierdo del triángulo, y solo entonces regresa
para explorar otras direcciones.

\begin{center}\rule{0.5\linewidth}{0.5pt}\end{center}

\subsection{Relleno Scanline}\label{relleno-scanline}

\subsubsection{Análisis
Matemático}\label{anuxe1lisis-matemuxe1tico-8}

El algoritmo Scanline optimiza el relleno al operar sobre \textbf{tramos
horizontales completos} en lugar de celdas individuales. Para cada
semilla \((s_x, s_y)\) extraída de la cola:

\begin{enumerate}
\def\labelenumi{\arabic{enumi}.}
\tightlist
\item
  \textbf{Búsqueda de límites:} escaneado lateral para encontrar
  \(x_{izq}\) y \(x_{der}\):
\end{enumerate}

\[x_{izq} = \min\{x \leq s_x : \text{matriz}[s_y][x] \neq \text{VACIA}\} + 1\]
\[x_{der} = \max\{x \geq s_x : \text{matriz}[s_y][x] \neq \text{VACIA}\} - 1\]

\begin{enumerate}
\def\labelenumi{\arabic{enumi}.}
\setcounter{enumi}{1}
\item
  \textbf{Relleno del tramo:} se pintan todas las celdas en
  \([x_{izq},\ x_{der}]\) de la fila \(s_y\).
\item
  \textbf{Propagación:} para cada celda pintada \((x, s_y)\), si la
  celda superior \((x, s_y-1)\) o inferior \((x, s_y+1)\) está vacía y
  aún no fue encolada, se agrega como nueva semilla.
\end{enumerate}

Este enfoque reduce el número de operaciones de encolado: en vez de
encolar cada celda individualmente, se encolan \textbf{puntos semilla de
fila} (uno por transición vertical), logrando un relleno mucho más
eficiente para figuras amplias.

\newpage

\subsubsection{Fragmento de Código}\label{fragmento-de-cuxf3digo-8}

\begin{lstlisting}
// RellenoScanlineController.cs — relleno por tramos horizontales
cola.Enqueue((semX, semY));
agregado[semY, semX] = true;

while (cola.Count > 0)
{
    var (sx, sy) = cola.Dequeue();
    if (matriz[sy, sx] != VACIA) continue;

    // Buscar límite izquierdo del tramo
    int xIzq = sx;
    while (xIzq > 0 && matriz[sy, xIzq - 1] == VACIA) xIzq--;

    // Buscar límite derecho del tramo
    int xDer = sx;
    while (xDer < COLS - 1 && matriz[sy, xDer + 1] == VACIA) xDer++;

    // Pintar el tramo completo y propagar semillas verticales
    for (int x = xIzq; x <= xDer; x++)
    {
        if (matriz[sy, x] != VACIA) continue;
        matriz[sy, x] = 2;
        pasos.Add(new PasoRelleno(x, sy, orden++));

        // Semilla superior
        if (sy > 0 && matriz[sy-1, x] == VACIA && !agregado[sy-1, x])
        { agregado[sy-1, x] = true; cola.Enqueue((x, sy-1)); }

        // Semilla inferior
        if (sy < ROWS-1 && matriz[sy+1, x] == VACIA && !agregado[sy+1, x])
        { agregado[sy+1, x] = true; cola.Enqueue((x, sy+1)); }
    }
}
\end{lstlisting}

\subsubsection{Resultado de Prueba}\label{resultado-de-prueba-8}

\textbf{Figura:} Pentágono centrado en \((11, 7)\) con radio \(5.5\)
celdas.

El algoritmo procesa la fila central primero (la que contiene la
semilla), rellena un tramo de \textasciitilde10 celdas en una sola
iteración de cola, y encola semillas para las filas \(6\) y \(8\). El
número de operaciones de encolado es proporcional al \textbf{número de
filas} de la figura (\textasciitilde9), no al número de celdas
(\textasciitilde70), lo que explica su ventaja en eficiencia para
figuras de gran ancho horizontal.


\section{Sistema de Coordenadas e Inversión del Eje
Y}\label{sistema-de-coordenadas-e-inversiuxf3n-del-eje-y}

Las API gráficas de sistemas operativos tradicionales (GDI+, DirectX 2D,
HTML Canvas) ubican el \textbf{origen \((0, 0)\) en la esquina superior
izquierda} de la ventana o panel, con el eje X creciendo hacia la
derecha y el \textbf{eje Y creciendo hacia abajo}. Esta convención es
opuesta a las coordenadas matemáticas estándar, donde el eje Y crece
hacia arriba.

Si se grafican directamente las coordenadas matemáticas sin
transformación, una línea con pendiente positiva (\(m > 0\)) aparecería
inclinada \textbf{hacia abajo} en la pantalla, y un círculo se vería
correcto en forma pero con el eje Y invertido respecto a las etiquetas.

Para corregir esto, el sistema aplica la siguiente transformación al
convertir coordenadas cartesianas \((x_c, y_c)\) a coordenadas de
pantalla \((p_x, p_y)\):

\[p_x = O_x + x_c \cdot e \qquad p_y = O_y - y_c \cdot e\]

Donde:

\begin{itemize}
\tightlist
\item \((O_x, O_y)\) es el \textbf{origen del plano} en píxeles de pantalla (centro del panel con desplazamiento por pan).
\item \(e\) es el \textbf{factor de escala} en píxeles por unidad matemática.
\end{itemize}

La inversión del signo en \(p_y\) (\(O_y - y_c \cdot e\)) corrige la
orientación del eje Y. Esta transformación está implementada en el
método \passthrough{\lstinline!CartesianoAPantalla!} de cada formulario:

\begin{lstlisting}
private PointF CartesianoAPantalla(float x, float y)
{
    PointF origen = ObtenerOrigen();
    return new PointF(
        origen.X + x * escala,   // X: dirección normal
        origen.Y - y * escala    // Y: invertido (GDI+ crece hacia abajo)
    );
}
\end{lstlisting}

La operación inversa, para convertir un clic del ratón en coordenadas
cartesianas, aplica la transformación opuesta:

\begin{lstlisting}
private PointF PantallaACartesiano(Point p)
{
    PointF origen = ObtenerOrigen();
    return new PointF(
        (p.X - origen.X) / escala,    // X normal
        (origen.Y - p.Y) / escala     // Y invertido
    );
}
\end{lstlisting}

El factor de escala inicial es de \(24\ \text{px/unidad}\), ajustable
mediante zoom desde \(10\) hasta \(72\ \text{px/unidad}\). El origen se
desplaza mediante clic-arrastre con el botón derecho, actualizando
\passthrough{\lstinline!desplazamientoPlano!} para implementar la
función de paneo.

\begin{center}\rule{0.5\linewidth}{0.5pt}\end{center}

\section{Documentación de Control de
Excepciones}\label{documentaciuxf3n-de-control-de-excepciones}

\subsection{Sanitización de Entradas
Numéricas}\label{sanitizaciuxf3n-de-entradas-numuxe9ricas}

Todos los campos de texto de coordenadas numéricos son procesados
mediante el método \passthrough{\lstinline!TryParse!} personalizado, que
intenta la conversión con dos culturas distintas para tolerar tanto la
coma decimal europea como el punto decimal anglosajón:

\begin{lstlisting}
private static bool TryParse(string texto, out float valor)
{
    // Intento 1: cultura del sistema operativo (ej. "3,14" en español)
    if (float.TryParse(texto, NumberStyles.Float,
        CultureInfo.CurrentCulture, out valor)) return true;

    // Intento 2: cultura invariante (ej. "3.14" en inglés)
    if (float.TryParse(texto, NumberStyles.Float,
        CultureInfo.InvariantCulture, out valor)) return true;

    valor = 0; return false;  // entrada inválida
}
\end{lstlisting}

Si cualquier campo falla la conversión, el sistema \textbf{no lanza una
excepción}, sino que retorna \passthrough{\lstinline!false!} y muestra
un \passthrough{\lstinline!MessageBox!} informativo, manteniendo el
estado previo de la aplicación intacto.

\subsection{Validación de Radio
Positivo}\label{validaciuxf3n-de-radio-positivo}

Para los algoritmos de cónicas, el radio debe ser estrictamente positivo
(\(r > 0\)). La validación se realiza en
\passthrough{\lstinline!LeerDatosDesdeTexto!}:

\begin{lstlisting}
private bool LeerDatosDesdeTexto()
{
    if (!TryParseInt(txtX1.Text, out int xc) ||
        !TryParseInt(txtY1.Text, out int yc) ||
        !TryParseInt(txtX2.Text, out int r)  ||
        r <= 0)               // radio debe ser > 0
        return false;
    // ...
}
\end{lstlisting}

Un radio de cero o negativo haría el algoritmo de Punto Medio
infinitamente iterativo (condición de parada \(x \leq y\) nunca se
cumpliría) y el algoritmo paramétrico generaría únicamente el punto
central.

\subsection{Validación de Semilla Dentro de la
Figura}\label{validaciuxf3n-de-semilla-dentro-de-la-figura}

La validación más sofisticada del sistema es la verificación de que la
semilla de relleno está \textbf{dentro} de la región cerrada y no en el
exterior. Se implementa como un BFS de pre-validación que busca un
camino libre de bordes hacia el límite de la matriz:

\begin{lstlisting}
private static bool EsDentroDelaFigura(byte[,] matriz, int semX, int semY)
{
    var visitado = new bool[ROWS, COLS];
    var cola = new Queue<(int X, int Y)>();
    cola.Enqueue((semX, semY));
    visitado[semY, semX] = true;

    while (cola.Count > 0)
    {
        var (x, y) = cola.Dequeue();

        // Si alcanza el borde de la matriz → semilla está fuera de la figura
        if (x == 0 || x == COLS-1 || y == 0 || y == ROWS-1)
            return false;

        foreach (var (nx, ny) in new[] { (x, y-1), (x+1,y), (x,y+1), (x-1,y) })
        {
            if (nx >= 0 && nx < COLS && ny >= 0 && ny < ROWS
                && matriz[ny, nx] == VACIA && !visitado[ny, nx])
            {
                visitado[ny, nx] = true;
                cola.Enqueue((nx, ny));
            }
        }
    }
    return true;  // la región explorada nunca alcanzó el borde → es interior
}
\end{lstlisting}

Esta validación tiene la misma complejidad \(O(n)\) que el relleno en
sí, pero se ejecuta sobre una copia de la matriz, dejando el estado
original intacto si la semilla es inválida.
\newpage

\subsection{Guardias de Nulidad en el Timer de
Animación}\label{guardias-de-nulidad-en-el-timer-de-animaciuxf3n}

El timer de animación verifica en cada tick que el estado del sistema
sea consistente antes de intentar acceder a datos:

\begin{lstlisting}
private void Timer_Tick(object? sender, EventArgs e)
{
    // Triple guardia: resultado, matriz y límite de pasos
    if (resultado == null || matriz == null || pasoActual >= resultado.Pasos.Count)
    {
        timer1.Stop();
        ActualizarEstado($"Relleno completo. Celdas pintadas: {resultado?.TotalCeldas}");
        return;
    }
    // Procesamiento seguro garantizado
    var paso = resultado.Pasos[pasoActual];
    if (paso.Y < ROWS && paso.X < COLS)
        matriz[paso.Y, paso.X] = 2;
    pasoActual++;
    pnlContenedor.Refresh();  // repintado síncrono para animación precisa
}
\end{lstlisting}

El uso de \passthrough{\lstinline!Refresh()!} en lugar de
\passthrough{\lstinline!Invalidate()!} garantiza que cada tick del timer
resulte en exactamente un frame visible, evitando que múltiples
invalidaciones se acumulen y produzcan saltos en la animación.

\begin{center}\rule{0.5\linewidth}{0.5pt}\end{center}

\section{Conclusiones Técnicas y Análisis de
Complejidad}\label{conclusiones-tuxe9cnicas-y-anuxe1lisis-de-complejidad}

\subsection{Tabla Resumen de
Algoritmos}\label{tabla-resumen-de-algoritmos}

{\def\LTcaptype{none} % do not increment counter
\begin{longtable}[]{@{}
  >{\raggedright\arraybackslash}p{(\linewidth - 8\tabcolsep) * \real{0.2051}}
  >{\raggedright\arraybackslash}p{(\linewidth - 8\tabcolsep) * \real{0.1453}}
  >{\raggedright\arraybackslash}p{(\linewidth - 8\tabcolsep) * \real{0.1795}}
  >{\raggedright\arraybackslash}p{(\linewidth - 8\tabcolsep) * \real{0.2650}}
  >{\raggedright\arraybackslash}p{(\linewidth - 8\tabcolsep) * \real{0.2051}}@{}}
\toprule\noalign{}
\begin{minipage}[b]{\linewidth}\raggedright
Algoritmo
\end{minipage} & \begin{minipage}[b]{\linewidth}\raggedright
Categoría
\end{minipage} & \begin{minipage}[b]{\linewidth}\raggedright
Estructura de datos
\end{minipage} & \begin{minipage}[b]{\linewidth}\raggedright
Técnica matemática
\end{minipage} & \begin{minipage}[b]{\linewidth}\raggedright
API principal
\end{minipage} \\
\midrule\noalign{}
\endhead
\bottomrule\noalign{}
\endlastfoot
DDA & Línea & Lista & Incremento fraccionario &
\passthrough{\lstinline!FillRectangle!} \\
Bresenham & Línea & Lista & Aritmética entera + error &
\passthrough{\lstinline!FillRectangle!} \\
Punto Medio Línea & Línea & Lista & Función implícita + decisión &
\passthrough{\lstinline!FillRectangle!} \\
Punto Medio Círculo & Cónica & Lista & Decisión incremental 8-fold &
\passthrough{\lstinline!FillRectangle!} \\
Paramétrico Círculo & Cónica & Lista + HashSet & Muestreo trigonométrico
& \passthrough{\lstinline!FillRectangle!} \\
Polar Círculo & Cónica & Lista + HashSet & Coordenadas polares &
\passthrough{\lstinline!FillRectangle!} \\
Relleno Semilla (BFS) & Relleno & Queue & Búsqueda en anchura &
\passthrough{\lstinline!FillRectangle!} \\
Relleno Pila (DFS) & Relleno & Stack & Búsqueda en profundidad &
\passthrough{\lstinline!FillRectangle!} \\
Scanline & Relleno & Queue & Tramos horizontales &
\passthrough{\lstinline!FillRectangle!} \\
\end{longtable}
}

\subsection{Análisis Comparativo de
Eficiencia}\label{anuxe1lisis-comparativo-de-eficiencia}

\subsubsection{Algoritmos de Líneas}\label{algoritmos-de-luxedneas}

{\def\LTcaptype{none} % do not increment counter
\begin{longtable}[]{@{}
  >{\raggedright\arraybackslash}p{(\linewidth - 6\tabcolsep) * \real{0.2338}}
  >{\raggedright\arraybackslash}p{(\linewidth - 6\tabcolsep) * \real{0.3766}}
  >{\raggedright\arraybackslash}p{(\linewidth - 6\tabcolsep) * \real{0.2208}}
  >{\raggedright\arraybackslash}p{(\linewidth - 6\tabcolsep) * \real{0.1688}}@{}}
\toprule\noalign{}
\begin{minipage}[b]{\linewidth}\raggedright
Algoritmo
\end{minipage} & \begin{minipage}[b]{\linewidth}\raggedright
Operaciones por píxel
\end{minipage} & \begin{minipage}[b]{\linewidth}\raggedright
Tipo aritmético
\end{minipage} & \begin{minipage}[b]{\linewidth}\raggedright
Complejidad
\end{minipage} \\
\midrule\noalign{}
\endhead
\bottomrule\noalign{}
\endlastfoot
DDA & 2 sumas, 2 redondeos & Punto flotante & \(O(k)\) \\
Bresenham & 2 sumas, 1 comparación & Entera & \(O(k)\) \\
Punto Medio & 1 suma, 1 comparación & Entera & \(O(k)\) \\
\end{longtable}
}

Donde \(k = \max(|\Delta x|, |\Delta y|)\). Aunque la complejidad
asintótica es equivalente, \textbf{Bresenham y Punto Medio son
significativamente más rápidos} en hardware que no cuenta con unidad de
punto flotante, dado que eliminan operaciones de redondeo y divisiones
en coma flotante. En hardware moderno con FPU, la diferencia es marginal
pero sigue siendo relevante en contextos de renderizado masivo.

\subsubsection{Algoritmos de
Circunferencia}\label{algoritmos-de-circunferencia}

{\def\LTcaptype{none} % do not increment counter
\begin{longtable}[]{@{}
  >{\raggedright\arraybackslash}p{(\linewidth - 6\tabcolsep) * \real{0.2118}}
  >{\raggedright\arraybackslash}p{(\linewidth - 6\tabcolsep) * \real{0.2471}}
  >{\raggedright\arraybackslash}p{(\linewidth - 6\tabcolsep) * \real{0.2000}}
  >{\raggedright\arraybackslash}p{(\linewidth - 6\tabcolsep) * \real{0.3412}}@{}}
\toprule\noalign{}
\begin{minipage}[b]{\linewidth}\raggedright
Algoritmo
\end{minipage} & \begin{minipage}[b]{\linewidth}\raggedright
Iteraciones totales
\end{minipage} & \begin{minipage}[b]{\linewidth}\raggedright
Tipo aritmético
\end{minipage} & \begin{minipage}[b]{\linewidth}\raggedright
Operaciones trigonométricas
\end{minipage} \\
\midrule\noalign{}
\endhead
\bottomrule\noalign{}
\endlastfoot
Punto Medio & \(\leq r/\sqrt{2}\) & Entera & Ninguna \\
Paramétrico & 360 & Punto flotante & 2 por iteración (cos, sin) \\
Polar & 360 & Punto flotante & 2 por iteración (cos, sin) \\
\end{longtable}
}

El \textbf{Punto Medio de Círculo} es el más eficiente: genera entre
\(\lfloor r/\sqrt{2}\rfloor + 1\) iteraciones (el primer octante), con
8-fold symmetry, totalizando \(O(r)\) operaciones enteras sin ninguna
función trigonométrica. El paramétrico y el polar son \(O(360)\)
independientemente del radio e involucran funciones transcendentes
costosas.

\subsubsection{Algoritmos de Relleno}\label{algoritmos-de-relleno}

{\def\LTcaptype{none} % do not increment counter
\begin{longtable}[]{@{}
  >{\raggedright\arraybackslash}p{(\linewidth - 6\tabcolsep) * \real{0.2024}}
  >{\raggedright\arraybackslash}p{(\linewidth - 6\tabcolsep) * \real{0.2857}}
  >{\raggedright\arraybackslash}p{(\linewidth - 6\tabcolsep) * \real{0.2619}}
  >{\raggedright\arraybackslash}p{(\linewidth - 6\tabcolsep) * \real{0.2500}}@{}}
\toprule\noalign{}
\begin{minipage}[b]{\linewidth}\raggedright
Algoritmo
\end{minipage} & \begin{minipage}[b]{\linewidth}\raggedright
Celdas encoladas
\end{minipage} & \begin{minipage}[b]{\linewidth}\raggedright
Patrón visual
\end{minipage} & \begin{minipage}[b]{\linewidth}\raggedright
Eficiencia relativa
\end{minipage} \\
\midrule\noalign{}
\endhead
\bottomrule\noalign{}
\endlastfoot
BFS (Semilla) & \(O(n)\) --- cada celda 1x & Uniforme concéntrico &
Base \\
DFS (Pila) & \(O(n)\) --- puede repetir & Profundidad variable & Similar
a BFS \\
Scanline & \(O(\text{filas})\) & Por filas & Mayor para figuras
amplias \\
\end{longtable}
}

El \textbf{Scanline} presenta la mejor eficiencia práctica para figuras
de geometría regular porque el número de semillas encoladas es
proporcional al número de \textbf{filas} de la figura, no al número de
celdas individuales. Para una figura de \(w\) columnas y \(h\) filas, el
Scanline encola \(O(h)\) semillas mientras BFS/DFS encolan hasta
\(O(w \times h)\) celdas, una mejora de un factor \(w\) en las
operaciones de gestión de la estructura de datos.

\subsection{\texorpdfstring{Decisión de API Gráfica:
\texttt{FillRectangle}
vs.~\texttt{DrawEllipse}}{8.3. Decisión de API Gráfica: FillRectangle vs.~DrawEllipse}}\label{decisiuxf3n-de-api-gruxe1fica-fillrectangle-vs.-drawellipse}

Todos los píxeles rasterizados se dibujan mediante
\passthrough{\lstinline!FillRectangle!} con un tamaño igual al factor de
escala, en vez de usar primitivas de alto nivel como
\passthrough{\lstinline!FillEllipse!} o
\passthrough{\lstinline!FillPolygon!}. Esta decisión tiene base en la
\textbf{semántica del algoritmo}: el objetivo es visualizar cada píxel
discreto como una celda cuadrada de la grilla, representando fielmente
el espacio de píxeles de un monitor real. Usar
\passthrough{\lstinline!FillEllipse!} o
\passthrough{\lstinline!FillPolygon!} para representar píxeles sería
semánticamente incorrecto, aunque visualmente más estético.

La figura \textbf{ideal} (la línea o circunferencia matemática continua)
sí utiliza las primitivas de alto nivel de GDI+
(\passthrough{\lstinline!DrawLine!},
\passthrough{\lstinline!DrawEllipse!}), que internamente aplican
antialiasing a nivel de subpíxel, creando la separación visual entre la
\emph{referencia matemática continua} y la \emph{aproximación discreta
rasterizada}.

\begin{center}\rule{0.5\linewidth}{0.5pt}\end{center}

\section{Repositorio}\label{repositorio}

El código fuente completo del proyecto, incluyendo el historial de
versiones, la estructura MVC completa y todos los algoritmos
documentados en este informe, está disponible en el repositorio de
GitHub del autor:

\textbf{Repositorio:}
\href{https://github.com/SrJCBM/Computacion-Grafica}{github.com/SrJCBM/Computacion-Grafica}

El repositorio contiene:

\begin{itemize}
\tightlist
\item Proyecto C\# SDK-style compilable con Visual Studio 2022 y .NET 10.0.
\item Estructura de carpetas MVC por algoritmo.
\item Historial de commits que refleja el desarrollo incremental.
\item Todos los formularios con interfaz interactiva paso a paso.
\end{itemize}

\begin{center}\rule{0.5\linewidth}{0.5pt}\end{center}

\end{document}
